\section{Preliminaries}
We model an agent interacting with an environment as a finite-horizon Markov Decision Process (MDP) $\mathcal{M} = (\mathcal{C}, \mathcal{S}, \mathcal{A}, P, R, H)$, where $\mathcal{C}$ denotes the space of contexts, $\mathcal{S}$ the state space, $\mathcal{A}$ the action space, $P : \mathcal{C} \times \mathcal{S} \times \mathcal{A} \rightarrow \Delta(\mathcal{S})$ the transition dynamics, $R : \mathcal{C} \times \mathcal{S} \times \mathcal{A} \rightarrow \mathbb{R}$ the reward function, and $H \in \mathbb{N}^+$ the horizon. At the beginning of each episode, a task prompt $x \in \mathcal{C}$ is sampled, and the agent begins in an initial state $s_1 \in \mathcal{S}$. At each timestep $t \in [1, H]$, the agent observes the current state $s_t$, selects an action $a_t \in \mathcal{A}$, and transitions to the next state $s_{t+1} \sim P(\cdot \mid x, s_t, a_t)$. In LLM-based agents, states correspond to prior interaction histories, and actions correspond to token sequences, such as natural language responses, code edits, and tool calls. A trajectory is defined as $\tau = (s_1, a_1, s_2, a_2, \dots, s_H, a_H)$. We assume access to a language model $\pi_\theta : \mathcal{C} \times \mathcal{S} \rightarrow \Delta(\mathcal{A})$, parameterized by $\theta$, from which actions are sampled autoregressively. A reward model assigns a scalar score to actions or trajectories. Conventional approaches rely on prompting LLMs to produce discrete scores in the language space. Formally, such reward models can be written as $R_{\text{LM}}(x, \tau) \in \{1, \dots, G\}$, where the score is the generated token.

